\exnumber{01}
\extitle{Project 3: Let's fair trade}
\exfiles{ex01/zip archive, named Name\_FirstName\_DeliverableName\_Date\_VersionNumber}
\exauthorize{Those named in the exercise}
\makeheaderfiles

Show a tested high-fidelity prototype.

\textbf{Estimated time:} 4 weeks

\textbf{Team size:} 3 people max

\textbf{Guidelines:}

\begin{itemize}
\item Test your prototypes with at least 5 people.
\item Make a presentation as if you were showing your work to Bruno. Make sure there are at least 5 people in the audience.
\item Upload the documents on Github, in a zip folder named "ProjectTitle\_Name\_FirstName", with all project deliverables as follows: "Name\_FirstName\_DeliverableName\_Date\_VersionNumber".  For example, the first deliverable can be named as "Doe\_John\_MidFiPrototype\_05242024\_V1".
\item Add a title within all written documents.
\item Pay attention to the global aspect of the deliverables. They must be clear, simple and easy to read. You can get inspiration on canva.com but don't overload your design with too much details!
\item Prepare a feedback form to send the audience. The form must include at least the questions listed below.
\item In both presentations, you are presenting to Barbara and her team.
\item Timing: 10 minutes for the presentation and 5 minutes for Q\&A.
\item You may use the presentation audits to structure your presentations.
\item Have fun!
\end{itemize}

\textbf{Tools:}

\begin{itemize}
\item Google Forms \url{https://docs.google.com/forms/u/0/}
\item Airtable \url{https://www.airtable.com/}
\item Typeform \url{https://www.typeform.com/}
\end{itemize}

\textbf{Tips:}

\begin{itemize}
\item Effective presentation skills \url{https://academic.oup.com/femsle/article/364/24/fnx235/4587905}
\end{itemize}

Intermediate deliverables:

\textbf{Quote:}

"Your mission in any presentation is to inform, educate, and inspire." Robert Ballard

\medskip

\textbf{Phase 1: Mid-project presentation audit}

Make sure:

\begin{itemize}
\item At least 5 people are attending the presentation
\item The presentation is about the whole UX process, including:
\begin{itemize}
\item The intial brief
\item A project planning, including the UX phases, the tools chosen, the timeline
\item User research findings from qualitative and quantitative research (persona, figures, etc.)
\item A problem statement to adapt the initial brief based on the research findings
\item An introduction of the concept that came out of the ideation
\item The animated mid-fidelity prototype (a demo-video)
\item The main feedback from usability testing
\item The next steps of the project
\end{itemize}
\item The presentation does not exceed 10 minutes
\item Feedback forms have been shared at the beginning of the presentation
\item You had fun!
\end{itemize}

\medskip

\textbf{Phase 2: Final presentation audit}

Make sure:

\begin{itemize}
\item At least 5 people are attending the presentation
\item The presentation is about the UI process, including:
\begin{itemize}
\item A quick reminder of the previous presentation
\item A project planning, including the UI phases, the tools chosen, the timeline
\item A moodboard
\item A styleguide (colors, branding, typography, some elements of the library)
\item The animated high-fidelity prototype (a demo-video)
\item The main feedback from usability and desirability testing
\item The next steps of the project
\end{itemize}
\item The presentation does not exceed 10 minutes
\item Feedback forms have been shared at the beginning of the presentation
\item You had fun!
\end{itemize}

\textbf{Questions to ask in the feedback forms}

\begin{itemize}
\item What do you remember of this presentation?
\item On a scale from 0 to 10, how clear was the speech?
\item Why?
\item On a scale from 0 to 10, how engaging was the presentation?
\item Why?
\item On a scale from 0 to 10, how coherent was the presentation?
\item Why?
\item On a scale from 0 to 10, how impactful were the visuals?
\item Why?
\item On a scale from 0 to 10, how clear was the problem?
\item Why ?
\item On a scale from 0 to 10, how well does the prototype answers the problem?
\item Why?
\item What recommendation would you like to make to improve this presentation?
\item Why?
\end{itemize}

